\chapter{随机最优控制、动态规划与 HJB 方程}
\lectureribbon{06}{HJB Equations, DPP and Stochastic Optimal Control I}{Andrzej Święch}
\begin{learninggoals}
\begin{itemize}
  \item 建立受控随机微分方程、成本泛函与值函数；
  \item 理解动态规划原理为何是“最优性的时间一致性”；
  \item 用短时间展开与 Itô 公式形式推导 HJB 方程；
  \item 认识 Nisio 半群如何把动态规划写成算子半群。
\end{itemize}
\end{learninggoals}

\section{受控扩散：决策与噪声共同推动状态}

在带滤过的概率空间 $(\Omega,\mathcal F,(\mathcal F_s)_{s\ge0},\Pp)$ 上，令 $W$ 为 $r$ 维 Brown 运动。控制取值于集合 $A$，容许控制 $\alpha_s$ 必须逐步可测，不能偷看未来噪声。状态方程为
\focusequation{受控随机微分方程}{%
\[
\dd X_s=b(s,X_s,\alpha_s)\dd s
+\sigma(s,X_s,\alpha_s)\dd W_s,
\qquad X_t=x.
\]}
漂移 $b$ 是可操纵的平均方向，扩散矩阵 $\sigma$ 决定随机扰动的形状。给定运行成本 $f$ 与终端成本 $g$，
\[
J(t,x;\alpha)=
\E\!\left[
\int_t^T f(s,X_s,\alpha_s)\dd s+g(X_T)
\right],
\qquad
V(t,x)=\inf_{\alpha\in\mathcal A_t}J(t,x;\alpha).
\]
$V(t,x)$ 把从 $(t,x)$ 出发的全部未来随机选择压缩成一个数。

\begin{center}
\begin{tikzpicture}[x=1cm,y=1cm]
  \draw[axis] (-4.4,-1.7)--(4.6,-1.7) node[right,font=\scriptsize\sffamily] {$s$};
  \fill[BlueAccent] (-4,-.2) circle (3.5pt);
  \foreach \k/\c/\yy in {1/BlueAccent/.8,2/PurpleAccent/.2,3/GreenAccent/-.55,4/YellowAccent/1.35,5/OrangeAccent/-1.1}{
    \draw[\c,thick] (-4,-.2)
      .. controls (-2.8,{-.1+.22*\k}) and (-1.6,{\yy-.45}) ..
      (-.6,\yy)
      .. controls (1.0,{\yy+.35-\k*.08}) and (2.6,{\yy-.25+\k*.06}) ..
      (4.2,{\yy+.15});
  }
  \node[font=\scriptsize\sffamily,text=BlueAccent] at (-3.45,.2) {同一初值};
  \node[conceptyellow] at (0,2.35) {控制改变整棵未来轨迹树的分布};
\end{tikzpicture}
\captionof{figure}{噪声产生许多可能路径；策略的目标不是指定一条轨迹，而是优化所有路径上的期望成本。}
\end{center}

\section{动态规划原理：把长问题切成两段}

对 $t\le\tau\le T$，动态规划原理（DPP）写成
\focusequation{Bellman 递归}{%
\[
V(t,x)=\inf_{\alpha\in\mathcal A_t}
\E\!\left[
\int_t^\tau f(s,X_s,\alpha_s)\dd s
+V(\tau,X_\tau)
\right].
\]}
前半段支付即时成本，后半段不再展开所有未来控制，而由 $V(\tau,X_\tau)$ 汇总。它要求最优策略的尾部在到达的随机状态上仍然最优，这就是\keyterm{时间一致性}。

\begin{center}
\begin{tikzpicture}[node distance=12mm and 10mm]
  \node[concept] (now) {现在\\$(t,x)$};
  \node[conceptgreen,right=of now] (tau) {中间状态\\$(\tau,X_\tau)$};
  \node[conceptyellow,right=of tau] (end) {终端\\$g(X_T)$};
  \draw[flow] (now)--node[above,font=\scriptsize\sffamily,text=MutedText]{$\int_t^\tau f\,\dd s$} (tau);
  \draw[warmflow] (tau)--node[above,font=\scriptsize\sffamily,text=MutedText]{$V(\tau,X_\tau)$} (end);
  \draw[softflow,bend right=28] (now) to node[below,font=\scriptsize\sffamily,text=BlueAccent]{一次优化 = 当前一步 + 最优余程} (end);
\end{tikzpicture}
\captionof{figure}{Bellman 原理把连续时间控制变成可无限细分的递归。}
\end{center}

\begin{pitfall}[DPP 不是普通的期望塔性质]
条件期望只能分解一个已给定策略的成本；DPP 还必须证明“前段控制”和“依随机中间状态选择的后段控制”能够可测地拼成容许控制。第八、九讲专门处理这一技术核心。
\end{pitfall}

\section{从 DPP 到 HJB：只看一个极短时间片}

定义固定控制 $a$ 下的扩散生成元
\[
\mathcal L^a\phi(t,x)
=b(t,x,a)\cdot D_x\phi(t,x)
+\frac12\Tr\!\left[\sigma\sigma^\top(t,x,a)D_x^2\phi(t,x)\right].
\]
在 $[t,t+h]$ 暂取常控制 $a$。若 $V$ 光滑，Itô 公式给出
\[
\E[V(t+h,X_{t+h})]-V(t,x)
=\E\int_t^{t+h}\bigl(\partial_sV+\mathcal L^aV\bigr)(s,X_s)\dd s.
\]
代入 DPP，除以 $h$ 并令 $h\downarrow0$，形式上得到
\focusequation{Hamilton--Jacobi--Bellman 方程}{%
\[
\partial_tV(t,x)+
\inf_{a\in A}\left\{
b(t,x,a)\cdot DV
+\frac12\Tr\!\left[\sigma\sigma^\top(t,x,a)D^2V\right]
+f(t,x,a)
\right\}=0,
\quad V(T,x)=g(x).
\]}

\begin{derivation}[每一项来自哪里]
\begin{itemize}
  \item $\partial_tV$：剩余时间缩短造成的价值变化；
  \item $b\cdot DV$：漂移沿价值曲面方向移动；
  \item $\tfrac12\Tr(\sigma\sigma^\top D^2V)$：Brown 运动的二次变差；
  \item $f$：当前瞬时成本；
  \item $\inf_a$：在每个 $(t,x)$ 上重新选择最佳动作。
\end{itemize}
\end{derivation}

\section{Nisio 半群：动态规划的算子形式}

对 $t\le s$ 和终端函数 $\varphi$，定义
\[
(\mathcal S_{t,s}\varphi)(x)
=\inf_{\alpha}
\E\!\left[
\int_t^s f(r,X_r,\alpha_r)\dd r+\varphi(X_s)
\right].
\]
DPP 等价于
\[
\mathcal S_{t,s}\circ\mathcal S_{s,r}
=\mathcal S_{t,r},\qquad t\le s\le r,
\]
并且 $V(t,\cdot)=\mathcal S_{t,T}g$。线性 Markov 半群传播期望，Nisio 半群则在传播中加入最小化，因此通常是非线性的。

\begin{intuition}[HJB 是半群的无穷小生成方程]
把 $s=t+h$，研究 $(\mathcal S_{t,t+h}\varphi-\varphi)/h$ 的极限，得到的正是 HJB 算子。于是 DPP 是有限时间尺度的递归，HJB 是它的微分形式。
\end{intuition}

\section{可算例：一维最小能量控制}

令 $\dd X_s=\alpha_s\dd s$，成本
\[
J(t,x;\alpha)=\int_t^T\frac12\alpha_s^2\dd s+\frac q2X_T^2,
\qquad q>0.
\]
HJB 为
\[
\partial_tV+\inf_a\{aV_x+\tfrac12a^2\}=0
\quad\Longrightarrow\quad
\partial_tV-\frac12(V_x)^2=0.
\]
作二次猜测 $V(t,x)=\tfrac12P(t)x^2$，得到
\[
P'(t)=P(t)^2,\qquad P(T)=q,
\qquad
P(t)=\frac{q}{1+q(T-t)}.
\]
最优反馈为
\[
\alpha^*(t,x)=-V_x(t,x)=-P(t)x.
\]
越靠近终端，$P(t)$ 越大，控制越积极地把状态拉回原点。

\begin{takeaway}
\begin{enumerate}
  \item 值函数把未来所有随机策略的最优期望成本压缩到当前状态；
  \item DPP 表达最优尾策略仍然最优；
  \item 短时间 DPP 加 Itô 公式产生 HJB；
  \item Nisio 半群是 DPP 的算子语言，HJB 是其无穷小生成方程。
\end{enumerate}
\end{takeaway}

\begin{exercisebox}
\begin{enumerate}
  \item 若 $\sigma$ 与控制无关，指出 HJB 中控制最小化只作用于哪些项。
  \item 对 $\dd X_s=\alpha_s\dd s+\eta\dd W_s$ 重做一维二次例子，观察扩散如何改变值函数常数项。
  \item 写出零运行成本时的 Nisio 半群，并与普通 Markov 半群比较。
\end{enumerate}
\end{exercisebox}
